\section{研究问题}

%%%%%%%%%%%%%%%
\begin{frame}{数据一致性问题}
  读操作语义问题:
  \begin{itemize}
	\setlength{\itemsep}{10pt}
	\item 共享数据环境: 返回\textcolor{red}{最新}写入的值
	\item 分布数据环境: ``最新''没有定义 \textcolor{red}{\small (副本)}
	\item \only<2->{\textcolor{red}{数据一致性}}问题: 在分布数据环境下, 读操作允许返回什么值?
  \end{itemize}
\end{frame}
%%%%%%%%%%%%%%%
\begin{frame}{数据一致性问题举例 (I)}
  \fig{width = 0.60\textwidth}{figures/data-inconsistency-comment-reordering.pdf}
  {社交网络中, 消息-评论乱序 \citeinbeamer{Lloyd}{CACM}{14}.}
\end{frame}
%%%%%%%%%%%%%%%
\begin{frame}{数据一致性问题举例 (II)}
  \begin{figure}[h!]
    \centering
    \begin{adjustbox}{max totalsize = {0.65\textwidth}{0.65\textheight}, center}
      %        File: data-inconsistency-rmw.tex
%     Created: Thu Oct 08 10:00 AM 2015 C
% Last Change: Thu Oct 08 10:00 AM 2015 C
%
\begin{tikzpicture}
  \uncover<1->{
  \begin{scope}
    % phone
    \node (phone) [] at (0,0) {\includegraphics[scale = 0.4]{figures/phone-icon.png}};
    % tablet
    \node (tablet) [below left = 4.0cm and 3.0cm of phone] {\includegraphics[scale = 
    0.55]{figures/tablet-icon.jpg}};
    % pc
    \node (pc) [below right = 4.0cm and 3.0cm of phone] {\includegraphics[scale = 
    0.60]{figures/pc-icon}};
  \end{scope}
  }

    % write from pc 
  \uncover<2->{
    \begin{scope}[<->, blue, line width = 6pt, font = \huge]
      % person at pc
      \only<2>{
      \node (person-pc) [above left = -1.0cm and -2.0cm of pc] {\includegraphics[scale = 
      0.30]{figures/person-icon}};
      }
      \draw [] (pc.west) to node [midway, sloped, above] {\textbf{1.} update file $f$} 
      (tablet.east);
      \draw [red, loosely dashed] (pc.north) to [out = 90, in = -20] node [midway, sloped, font = 
      \Huge, scale = 2]{$\times$} (phone.east);
      \draw (pc) to [loop, out = 60, in = -10, looseness = 3] node [midway, above, sloped] 
      {\textbf{1.} update file $f$} (pc);
    \end{scope}
   }


    % read from phone
   \uncover<3->{
    \begin{scope}[<->, brown, line width = 6pt, font = \huge]

      % person at phone
      \node (person-phone) [above left = -1.0cm and -2.0cm of phone] {\includegraphics[scale = 
      0.30]{figures/person-icon}};

      \draw (phone) to [loop, out = 180, in = -100, looseness = 3] node [midway, below, sloped] 
      {\textbf{2.} read file $f$} (phone);   
    \end{scope}

    % update lost
    \node (inconsistency) [font = \huge, right = of person-phone, red] {\bf Update Lost!};
  }
\end{tikzpicture}

    \end{adjustbox}
    \caption{多设备文件共享时, 更新丢失 ($\#N = 3, \#W = 2, \#R = 1$).}
  \end{figure}
\end{frame}
%%%%%%%%%%%%%%%
\begin{frame}{分布共享数据服务}
  \begin{columns}
	\column{0.50\textwidth}
	  \todo{图: 分布共享数据服务 (简化图)}
	  (数据)一致性模型是核心概念.
	\column{0.50\textwidth}
	  \begin{itemize}[<+->]
		\setlength\itemsep{10pt}
		\item 简化上层应用的开发
		\item 提供共享数据的抽象
		\item 屏蔽底层数据的分布性
	  \end{itemize}
  \end{columns}
\end{frame}
%%%%%%%%%%%%%%%
\begin{frame}{数据一致性模型}
\end{frame}
%%%%%%%%%%%%%%%
% \begin{frame}{数据一致性问题}
%   数据一致与否是相对于应用逻辑而言的:
%   \pause
%   \begin{itemize}[<+->]
% 	\setlength\itemsep{8pt}
% 	\item 数据一致性模型多样
% 	\item 数据一致性模型有强弱之分
% 	\item 应用规约数据一致性需求
% 	\item 数据不一致导致应用异常(anomalies)
%   \end{itemize}
% \end{frame}
%%%%%%%%%%%%%%%
\begin{frame}{数据一致性问题研究的历史阶段}
  \fignocaption{width = 0.30\textwidth}{figures/taste-of-past.jpg}

  \todo{图: ps}
\end{frame}
%%%%%%%%%%%%%%%
\begin{frame}{数据一致性问题研究的历史阶段}
  \fignocaption{width = 0.90\textwidth}{figures/consistency-model-history.pdf}
\end{frame}
%%%%%%%%%%%%%%%
\begin{frame}{数据一致性问题研究的发展趋势}
  \begin{center}
	云计算凸显应用价值观
  \end{center}

  \fignocaption{width = 0.40\textwidth}{figures/sla-tools.jpg}
\end{frame}
%%%%%%%%%%%%%%%
\begin{frame}{数据一致性问题研究的发展趋势 (I)}
  \begin{description}
	\setlength{\itemsep}{20pt}
	\item[精细化:] 从二元到连续谱
	  \begin{columns}
		\column{0.50\textwidth}
		  \fignocaption{scale = 0.06}{figures/coin-flip.jpg}
		\column{0.50\textwidth}
		  \fignocaption{scale = 0.05}{figures/prism-spectrum.jpg}
	  \end{columns}
	\item[可度量:] think \textcolor{red}{\it probabilistically} \citeinbeamer{Brewer}{}{}
	  \fignocaption{width = 0.25\textwidth}{figures/tape-measure.jpg}
	  \begin{center}
		量化系统执行, 后验系统对一致性的满足程度
	  \end{center}
  \end{description}
\end{frame}
%%%%%%%%%%%%%%%
\begin{frame}{数据一致性问题研究的发展趋势 (II)}
  \begin{description}
	\setlength{\itemsep}{20pt}
	\item[多样化:] 从单一到???
	  \begin{itemize}
		\item 一致性族: causality; read-your-writes {\scriptsize (RYW)}
		\item 参数调节: 提供``有限度''的不一致 \citeinbeamer{Yu}{TOCS}{02}
	  \end{itemize}
	  \fignocaption{width = 0.30\textwidth}{figures/one-size-fit-all-table.jpg}
	\item[可调节:] 混合使用, 运行时可变
	  % \fignocaption{width = 0.20\textwidth}{figures/rudder.jpg}
  \end{description}
\end{frame}
%%%%%%%%%%%%%%%
\begin{frame}{我们的工作}
  
\end{frame}
%%%%%%%%%%%%%%%
% \begin{frame}{数据一致性问题研究的发展趋势}
%   \begin{enumerate}
%     \setlength{\itemsep}{15pt}
% 	\item 精细化, 可度量 \citeinbeamer{Bailis}{VLDB}{12}
% 	  \vspace{5pt}
% 	  \uncover<2->{
% 		\begin{description}
% 		  \setlength{\itemsep}{5pt}
% 		  \item[精细化:] ``在大多数情况下, 访问到一致数据''
% 		  \item[可度量:] 量化系统执行, 后验系统对一致性的满足程度 
% 		\end{description}
% 	  }
% 	  \only<2->{
% 		\begin{textblock*}{6cm}(6.0cm, 1.5cm)
% 		  \fignocaption{width = 0.60\textwidth}{figures/tape-measure.jpg}
% 		\end{textblock*}
% 	  }
%     \item 多样化, 可定制 \citeinbeamer{Terry}{CACM}{13}
% 	  \vspace{5pt}
% 	  \uncover<3->{
% 		\begin{description}
% 		  \setlength{\itemsep}{5pt}
% 		  \item[多样化:] 
% 			\begin{itemize}
% 			  \item 一致性族: causality; read-your-writes {\scriptsize (RYW)}
% 			  \item 参数调节: 提供``有限度''的不一致 \citeinbeamer{Yu}{TOCS}{02}
% 			\end{itemize}
% 		  \item[可定制:] 混合使用, 运行时可变
% 		\end{description}
% 	  }
% 	  \only<3->{
% 		\begin{textblock*}{6cm}(6.0cm, 6.5cm)
% 		  \fignocaption{width = 0.35\textwidth}{figures/rudder.jpg}
% 		\end{textblock*}
% 	  }
%   \end{enumerate}
% \end{frame}
%%%%%%%%%%%%%%%
%%%%%%%%%%%%%%%
%%%%%%%%%%%%%%%
%%%%%%%%%%%%%%%
